\section{Heat flow and the complementary physical space}\label{sec:heat}

The second manuscript of the latest edition, \texttt{HEAT\_AND\_COMPLEMENT.md} (H1--H34; Appendix~B of the fifth-reference reader), uses the control of the vacuum from Section~\ref{sec:gap} for two further purposes: box-independent bounds on the heat correlations of plaquettes, and bounds on how much the plaquette family mixes with the rest of the physical space. This section states both, with proofs of the steps that carry the constants.

\subsection{Objects}

Let $A=H_L-E_0$ on the centred physical space and $\rho=\psi^2$. Every statement below holds for every $L\ge2$ and $a>0$.
\begin{itemize}
\item \emph{Centred plaquette states and their heat correlations}:
\[
r_p=(W_p-\langle W_p\rangle_\rho)\psi,\qquad \widehat C_{pq}(\tau;\xi)=\langle r_p,e^{-\tau A/\kappa}r_q\rangle,\qquad\tau=\kappa t.
\]
\item \emph{Row norm}: $\|B\|_{\mathrm{row}}=\max_p\sum_q|B_{pq}|$.
\item \emph{Taylor coefficients}: $C_0,C_2,C_4$ are the coefficients of $\xi^0,\xi^2,\xi^4$ in $\widehat C$; the workbench computed them in its 21 September heat edition.
\item \emph{The transformed generator}: in the coordinates $f\mapsto\psi f$, Lemma~\ref{lem:gst} turns $A/\kappa$ into $K-D_v$ with $D_vh=2\Gamma(v,h)$.
\item \emph{The Fourier-algebra norm} is $\|\cdot\|_X$ of Section~\ref{sec:gap}, and $\chi=4\sum_it_i\xi^i+\frac23w_*$ is the number from the proof of Proposition~\ref{prop:return}, so $3(1-\chi)=d_5(\xi)$ and $\|D_vh\|_X\le\chi\|Kh\|_X$.
\end{itemize}

\subsection{A heat remainder independent of the box}

\begin{theorem}[heat remainder; H13]\label{thm:heat}
For every $L\ge2$, every $\tau\ge0$ and every $|\xi|<1/55$,
\[
\bigl\|\widehat C(\tau;\xi)-C_0(\tau)-\xi^2C_2(\tau)-\xi^4C_4(\tau)\bigr\|_{\mathrm{row}}\le\frac{67896}{169}\,e^{-13\tau/8}\,\frac{(55|\xi|)^6}{1-(55|\xi|)^2}.
\]
\end{theorem}
\status{\textsf{Audited}, given the analytic core of Theorem~\ref{thm:ym01}, on which it rests; it is therefore conditional on the same order-3 to order-5 inputs. Every constant is re-derived below. The semigroup construction (H5--H7) and the analytic continuation to complex $\xi$ were read and are standard; the evenness and the Cauchy step are proved below. The inherited matrices $C_0,C_2,C_4$ are not needed for the truth of the bound, only for its use. The earlier edition's bound, on the circle $3/256$ with prefactor $512$ and rate $3/2$, is superseded.}

The proof is in four steps; the constants are collected in Proposition~\ref{prop:heatconst}.

\emph{(i) The mean and the semigroup} (H1--H7). Put $T=Q_HD_vK^{-1}$ and $p=P_HD_vK^{-1}$ on the nonconstant physical coefficients. Since the norm $\|\cdot\|_X$ adds the scalar and the nonconstant parts, $|py|+\|Ty\|_X\le\chi\|y\|_X$. So $S=(I-T)^{-1}$ exists with $\|S\|\le1/(1-\chi)$, and $\mu(F)=P_HF+pSQ_HF$ satisfies $|(\mu-P_H)F|\le\chi\|Q_HF\|_X/(1-\chi)$.
For $L_0=(I-T)K$ and $h>0$, coefficientwise $\|(I+hK)f\|_X=\|f\|_X+h\|Kf\|_X$. Hence
\[
\|(I+hL_0)f\|_X\ge\|f\|_X+h(1-\chi)\|Kf\|_X\ge(1+3h(1-\chi))\|f\|_X,
\]
using $c_j\ge3$ on nonconstant physical blocks (Lemma~\ref{lem:casimir}). With the range condition (H5), the Lumer--Phillips theorem \cite{LumerPhillips} gives a semigroup $E(\tau)=e^{-\tau L_0}$ with $\|E(\tau)\|\le e^{-3(1-\chi)\tau}$.

\emph{(ii) A formula for the correlations.} Let $\mathcal T(\tau)$ be the transformed heat semigroup, so that $\widehat C_{pq}(\tau)=\mu[(\mathcal T(\tau)F_p)W_q]$ with $F_p=W_p-\mu W_p$; here $\mu[\mathcal T(\tau)F_p]=\mu[F_p]=0$ removes the second mean. Let $h_\tau=K^{-1}SE(\tau)W_p$ and $G_z=\sum_qz_qW_q$ with $|z_q|\le1$. Since $W_p$ has Haar mean $0$, the Poisson equation (H3) gives $(K-D_v)h_\tau=E(\tau)W_p-\mu(E(\tau)W_p)=\mathcal T(\tau)F_p$. For any $h$ and $G$,
\[
\int e^{2v}\bigl(Kh-2\Gamma(v,h)\bigr)G\,dU=\int e^{2v}\,\Gamma(h,G)\,dU,
\]
by moving one derivative $X_{e,a}$ onto $e^{2v}G$. Dividing by $\int e^{2v}$ gives $\sum_qz_q\widehat C_{pq}(\tau)=\mu\Gamma(h_\tau,G_z)$, and $\|Kh_\tau\|_X\le\|S\|\,\|E(\tau)\|\,\|W_p\|_X\le\frac8{1-\chi}e^{-3(1-\chi)\tau}$.

\begin{proposition}[the constants of the heat bound; H9--H12]\label{prop:heatconst}
In this setting:
\begin{enumerate}[label=(\alph*)]
\item $\|\Gamma(h,G_z)\|_X\le32\|Kh\|_X$;
\item $|P_H\Gamma(h,G_z)|\le\frac18\|Kh\|_X$;
\item $\|\widehat C(\tau)\|_{\mathrm{row}}\le\dfrac{1+255\chi}{(1-\chi)^2}\,e^{-3(1-\chi)\tau}$;
\item at $|\xi|=1/55$ one has $1/55<\alpha_5$ and $d_5(1/55)=1.637\ldots>13/8$, so $\chi<11/24$, the prefactor is at most $\frac{1+255\cdot11/24}{(13/24)^2}=\frac{67896}{169}$, and the rate is at least $\frac{13}8$.
\end{enumerate}
\end{proposition}
\status{Proved here. The workbench's (H9)--(H12) state the same; the factor $255$ is derived here. The inequalities of (d) are checked in exact rational arithmetic (script \note{ymr\_heat\_constants.py}).}
\begin{proof}
(a) By (F10), $\|\Gamma(h,G)\|_X\le2t(G)\|Kh\|_X$. The function $W_q$ has one Fourier block, with spin $\frac12$ on the four links of $q$, and its coefficient matrix has trace norm $2^{4-1}=8$ by (F12): four fundamental occurrences, two cyclic sign changes. A link lies on at most four plaquettes, so $t(G_z)\le4\cdot\frac12\cdot8=16$.
(b) Integration by parts gives $P_H\Gamma(h,G_z)=\int h\,KG_z\,dU=3\sum_qz_q\int hW_q\,dU$, since each $W_q$ has Casimir $4\cdot\frac34=3$. The gauge-invariant functions whose Fourier support is spin $\frac12$ on the four links of $q$ and spin $0$ elsewhere form the line $\C W_q$: at each corner two spin-$\frac12$ links meet, and $\frac12\otimes\frac12$ contains the trivial representation once. Write $a_q$ for the component of $h$ on this line. By Haar orthogonality and $\int W_q^2=1$, $\int hW_q=a_q$. The block of $a_qW_q$ contributes $3\cdot8|a_q|$ to $\|Kh\|_X$, so $|P_H\Gamma|\le3\sum|a_q|\le\frac18\|Kh\|_X$.
(c) $\mu\Gamma=P_H\Gamma+(\mu-P_H)Q_H\Gamma$, so by (b), step (i) and (a),
\[
|\mu\Gamma(h_\tau,G_z)|\le\Bigl(\frac18+\frac{32\chi}{1-\chi}\Bigr)\|Kh_\tau\|_X\le\Bigl(\frac18+\frac{32\chi}{1-\chi}\Bigr)\frac{8e^{-3(1-\chi)\tau}}{1-\chi}=\frac{1+255\chi}{(1-\chi)^2}e^{-3(1-\chi)\tau}.
\]
The row norm is the maximum over the phases $z_q$ of the left side.
(d) The value $d_5(1/55)$ is computed from the rationals of (F26), and $d_5>13/8$ is equivalent to $\chi<11/24$ since $d_5=3(1-\chi)$. The prefactor increases with $\chi$, and $3(1-\chi)>3\cdot\frac{13}{24}=\frac{13}8$.
\end{proof}

\emph{(iii) Evenness.} The map $U_{(n,i)}\mapsto(-1)^{\sum_{k<i}n_k}U_{(n,i)}$ multiplies each link by a central element. It preserves the Haar measure, $K$ and the gauge action, and it sends every $W_p$ to $-W_p$: for $p=(n;i,j)$ with $i<j$ the four signs multiply to $-1$. It therefore conjugates $H_L(\xi)$ to $H_L(-\xi)$ plus a constant, and maps each $r_p$ to $-r_p$ at the opposite coupling. So $\widehat C(\tau;-\xi)=\widehat C(\tau;\xi)$, and only even powers of $\xi$ occur. (Checked on all $240$ faces of the box $L=2$.)

\emph{(iv) Cauchy's estimate.} By (c) and (d), on the circle $|\zeta|=1/55$ every phase-weighted row sum of $\widehat C(\tau;\zeta)$ is at most $Ce^{-13\tau/8}$ with $C=67896/169$. So the coefficient of $\zeta^{2n}$ has row norm at most $C e^{-13\tau/8}55^{2n}$, and summing over $n\ge3$ gives the stated tail $\sum_{n\ge3}(55|\xi|)^{2n}=(55|\xi|)^6/(1-(55|\xi|)^2)$. This proves Theorem~\ref{thm:heat}.

\emph{Other radii} (\textsf{Generalised here}). Nothing in steps (i)--(iv) uses $R=1/55$ except the numbers of Proposition~\ref{prop:heatconst}(d). For every $0<R<\alpha_5$ the same proof gives the bound of Theorem~\ref{thm:heat} with $55|\xi|$ replaced by $|\xi|/R$, the prefactor replaced by $C(R)=(1+255\chi_R)/(1-\chi_R)^2$ and the rate $13/8$ replaced by $d_5(R)=3(1-\chi_R)$, where $\chi_R=1-d_5(R)/3$. The radius can therefore approach $\alpha_5\approx1/54.27$. For example $R=0.0184$ gives $C(R)=443.37\ldots$ and rate $1.5746\ldots$. At $R=1/55$ itself the exact $\chi_R=0.45412\ldots$ gives the prefactor $391.98\ldots$ and the rate $1.6376\ldots$; the workbench's rational rounding $\chi<11/24$ costs little. (Computed from the rationals of (F26); these are floating-point values, not rational certificates.)

\subsection{Relaxation into the whole complementary space}\label{ssec:complement}

Let $R$ be the map with columns $r_p$ and $\Phi=\kappa A^{-1}R$. The original Grams and the energy of the family are
\[
G_0=R^*R,\quad G_1=R^*\Phi,\quad G_2=\Phi^*\Phi,\quad E=\kappa G_1,\quad K_o=\kappa^{-1}R^*AR.
\]
Put $P=\Phi G_2^{-1}\Phi^*$, $Q=I-P$, $B_\Phi=QR$ and $D_Q=QAQ$ on $Q\mathcal H_0$. Allowing the family $\Phi c$ to relax into every $h\in Q\mathcal H_0$ gives the relaxed energy and Gram
\[
E_{\mathrm{full}}=E-\kappa^2B_\Phi^*D_Q^{-1}B_\Phi,\qquad G_{\mathrm{full}}=G_2+\kappa^2B_\Phi^*D_Q^{-2}B_\Phi .
\]

\begin{proposition}[relaxation bounds; H20--H29]\label{prop:relax}
Suppose that, for numbers $w_0,w_1,w_2,w_K,\beta\ge0$ and $d_{\mathrm{ph}}>0$,
\[
\|G_k-3^{-k}I\|\le w_k\ (k=0,1,2),\qquad\|K_o-3I\|\le w_K,\qquad 0\preceq J=G_0-6G_1+9G_2\preceq\beta I,\qquad A\ge\kappa d_{\mathrm{ph}}.
\]
Put $l_k=3^{-k}-w_k>0$, $u_k=3^{-k}+w_k$ and $u_K=3+w_K$. Then:
\begin{enumerate}[label=(\alph*)]
\item the fraction of each $\Phi c$ seen by the plaquette states, $\|P_R\Phi c\|^2/\|\Phi c\|^2$, is at least $l_1^2/(u_0u_2)$, where $P_R$ is the orthogonal projection onto the span of the $r_p$;
\item $(1-\eta_E)E\preceq E_{\mathrm{full}}\preceq E$ and $G_2\preceq G_{\mathrm{full}}\preceq(1+\eta_G)G_2$, with $\eta_E=\beta/(d_{\mathrm{ph}}l_1)$ and $\eta_G=\beta/(d_{\mathrm{ph}}^2l_2)$;
\item the energy of $(I-P_R)\Phi c$ is at most $(u_1u_K/l_0^2-1)$ times $c^*Ec$.
\end{enumerate}
\end{proposition}
\status{Proved here; the workbench's (H21), (H23), (H26)--(H29). The domain statement for $D_Q$ (H25) is the workbench's and was read.}
\begin{proof}
(a) $\|P_R\Phi c\|^2=c^*G_1G_0^{-1}G_1c\ge u_0^{-1}c^*G_1^2c\ge(l_1^2/u_0)|c|^2$ and $\|\Phi c\|^2=c^*G_2c\le u_2|c|^2$.
(b) $B_\Phi^*B_\Phi=G_0-G_1G_2^{-1}G_1$ is the least value of $(R-\Phi X)^*(R-\Phi X)$ over coefficient matrices $X$, and $X=3I$ gives $J$; so $B_\Phi^*B_\Phi\preceq J\preceq\beta I$. Since $D_Q\succeq\kappa d_{\mathrm{ph}}$, $\kappa^2B_\Phi^*D_Q^{-1}B_\Phi\preceq(\kappa/d_{\mathrm{ph}})\beta I\preceq(\beta/(d_{\mathrm{ph}}l_1))\kappa G_1$, and $\kappa^2B_\Phi^*D_Q^{-2}B_\Phi\preceq(\beta/d_{\mathrm{ph}}^2)I\preceq(\beta/(d_{\mathrm{ph}}^2l_2))G_2$.
(c) Since $A\Phi=\kappa R$ and $P_RR=R$, the mixed energy is $q_A(\Phi c,P_R\Phi c)=\kappa\langle Rc,\Phi c\rangle=\kappa c^*G_1c=q_A(\Phi c)$, and $P_R\Phi c=RG_0^{-1}G_1c$ has energy $\kappa c^*G_1G_0^{-1}K_oG_0^{-1}G_1c$. Expanding $q_A((I-P_R)\Phi c)$ gives $\kappa c^*(G_1G_0^{-1}K_oG_0^{-1}G_1-G_1)c$, the workbench's (H29). Now $G_1G_0^{-1}K_oG_0^{-1}G_1\preceq(u_K/l_0^2)G_1^2\preceq(u_Ku_1/l_0^2)G_1$.
\end{proof}

The workbench obtains the numbers $w_k$, $w_K$ and $\beta$ from Theorem~\ref{thm:heat}. The time-integral formula $G_k=\int_0^\infty\tau^{k-1}\widehat C(\tau)\,d\tau/(k-1)!$ ($k\ge1$) gives the remainder factors $C/d^k$ with $d=13/8$. For $J=\widehat C(0)+\int_0^\infty(9\tau-6)\widehat C(\tau)\,d\tau$ the signed weight is kept: $\int_0^\infty|9\tau-6|e^{-d\tau}d\tau=\frac6d-\frac9{d^2}+\frac{18}{d^2}e^{-2d/3}$, and with $e^{-13/12}<\frac{339}{1000}$ the prefactor is $C_J=5156090136/3570125$ (H17--H18). The degree-2 and degree-4 parts are the workbench's inherited row bounds (the table after H18), which come from 199 supports and 559 marked contributions of its earlier volume edition.

\begin{proposition}[the benchmark table; H5]\label{prop:bench}
From the stated inherited row bounds and the constants $C$, $C_J$ and $606$, the widths and Proposition~\ref{prop:relax} give, for every $L\ge2$ and $a>0$:

\begin{center}\small
\begin{tabular}{@{}lccc@{}}
\toprule
Coupling & Observed fraction & Energy loss $\eta_E$ & Gram increase $\eta_G$\\
\midrule
$g^2\ge10$ & $>0.9778$ & $<0.01057$ & $<0.01123$\\
$g^2\ge12$ & $>0.9975$ & $<0.00115$ & $<0.00120$\\
$g^2\ge25/2$ & $>0.9984$ & $<0.00071$ & $<0.00073$\\
$g^2\ge13$ & $>0.99904$ & $<0.000436$ & $<0.000451$\\
$g^2\ge16$ & $>0.99991$ & $<0.0000357$ & $<0.0000364$\\
\bottomrule
\end{tabular}
\end{center}
In particular the workbench's table holds: at $g^2\ge13$ the energy loss and the Gram increase are below $1/2000$, and at $g^2\ge16$ below $1/25000$.
\end{proposition}
\status{\textsf{Reproduced}: recomputed in exact rational arithmetic from the stated inputs (script \note{ymr\_heat\_constants.py}); at $g^2=13$, with the workbench's $d_{\mathrm{ph}}=29047/10000$, the three values agree with its diagnostics to eleven digits. At the other couplings the script uses $d_{\mathrm{ph}}$ equal to $d_5(\xi)$ rounded down to four digits. \textsf{Conditional} on the inherited degree-2 and degree-4 row bounds and on the kinetic-row constants $30$ and $48$ of (H14), which were not re-derived, and on the inputs of Theorem~\ref{thm:ym01}. The widths are monotone in $\xi$ and $d_5$ decreases in $\xi$, so each row holds for all larger $g^2$, as the workbench says.}

\subsection{Heat horizons, support quotients and the spatial limit}

\begin{itemize}
\item \emph{Finite heat time} (H33; the Gram comparison is the workbench's, the choice of $T$ is checked here). Replacing $\Phi$ by $\Phi_T=(I-e^{-TA})\Phi$ changes the Grams by factors between $(1-\varepsilon)^2$ and $1$, with $\varepsilon=e^{-\kappa d_{\mathrm{ph}}T}$. For a signed sum of $\log\det$ terms of total rank $B$ the error is at most $2B(-\log(1-\varepsilon))\le2B\varepsilon/(1-\varepsilon)$, and the choice $T=(\kappa d_{\mathrm{ph}})^{-1}\log(1+2B/\eta)$ makes $\varepsilon/(1-\varepsilon)=\eta/(2B)$, so the error is at most $\eta$.
\item \emph{Support quotients} (H31--H32; located). An exact isomorphism between quotients of inverse-energy spans and kernels, and the attained minimum of the family energy over changes of support.
\item \emph{The spatial limit at fixed spacing} (H34; located, an outline). Heat coefficients through degree $N$ agree on two boxes containing the same $N$-neighbourhood of the marks, so Theorem~\ref{thm:heat} makes the correlations converge as $L\to\infty$; the workbench then builds the infinite-volume semigroup, its invariant state and the relaxation bounds. This is a limit at fixed $a$ and $g$, not a continuum limit, as the workbench says.
\end{itemize}
